%===================================== CHAP 1 =================================

\chapter{Introduction}

\section{Background and motivation}
Unmanned aerial vehicles have in seen an increase usage in the civilian sector on land, where they can be used at a cheaper price then manned aircraft. UAV are well suited for tasks like inspections, aerial photography, environment surveillance and search and rescue. Today UAV are mostly operated over land, but in the future this will include the out at sea as well. This will give some challenges that must be overcome. One of these challenges is that the UAV should be able to perform a autonomous landing.

A UAV can increase performance in many maritime operation where today only other manned aircraft or satellites are the only solution. In the maritime sector they can be used in iceberg management, monitoring of oil spills, search and rescue and maritime traffic monitoring. To enable a safe UAV operation at sea there must be a system in place to ensure a safe landing.

There exist landing system that can guide the UAV towards a net, but they are expensive and restricted to a few UAVs. A pilot could always land the UAV, but it would be better if the UAV could hit the net by it self. In order to make the UAV able to perform a automatic landing the minimum requirement is that it know where it is at all time. This put a requirement of the position sensor. The position system need to combine low cost with high position accuracy. This thesis will test a new generation of GNSS receiver, and use GPS to find the position to the UAV. The demand on the accuracy is that the error must be in decimetres to ensure safe landing in the net.

A automatic landing system must include a path planner, guidance system and a accurate position estimation system, in addition to the low level control system in the UAV. A path planner and guidance system were created in a master thesis at the uavlab by Marcus Frølich. The gudiance system is currently under further development in two project thesis. A position system was created by Bjørn Spockeli, but it concluded that the GNSS receiver used were insufficient to perform automatic landing. This thesis will continue the research done by Spockeli, and introduce a new GNSS receiver that will be used with the open source program, rtklib. The motivation is to have a system with accurate local position estimate, such that in the future a landing can be performed automatic.
The automatic landing system will use RTK GPS for position estimation. 
Motivation for work: Autonomos landing. Required accurate posisioning system, a feasable path, tuned guidence system, robust lav nivå kontroll. Include a system figure can give the reader information on how the autolanding system should work, and why this thesis is relevant.
\section{Previous work}
Citation checking \citep{Spockeli,Froelich} or \cite{Spockeli,Froelich} or \citet{Spockeli,Froelich}

GNSS navigation has been
As part of the NTNU MSc thesis of \citep{Skulstad&Syversen} they successfully demonstrated a net landing with a fixed wing UAV using a low-cost RTK GNSS system. More about how they did it.

An other NTNU MSc thesis by \citep{Spockeli} simulated a net landing with a fixed wing UAV. Here the net was allowed to be dynamical, and the system was design with landing on a ship in mind. The system used artificial neural network to estimate the position of the net.

As part of the Stellenbosch University MSc thesis \citep{smit2013autonomous} demostraded a autonomous landing on a airfield using decoupled linear controllers and DGPS. 

An other vision \citep{kim2013fully}

The uavlab at NTNU has in the last year studied how to perform a autonomous landing with a fixed wing UAV. Automatic landing of a UAV in a net using low-cost single frequency gps is described in \citep{Skulstad&Syversen}, and a path generation system that is integraded with neptus is described in \citep{Froelich}. Research on RTK GPS intergration in Dune is described in \citep{Spockeli}.

Work done on ambiguity
\cite{Ambiguity:Estimation,LAMBDA:METHOD,LAMBDAMETHOD,GeodeticBaselines} 
Work done with rtk gps
\citep{Low-costRTK}  \citep{3D-RTK}

There has been done work on autonomous landing system using a vision aided system. The work done in \citep{williams2012intelligent} describe a intelligent vision aided landing system that detect and generate landing waypoints for a unsurveyed airfield. Other work 

Commercial landing system:  \citep{scanealge}

\section{Goal of thesis}
The goal of this thesis is to test the performance of ublox-LEA M8T GNSS receiver in a real time differential \gls{dgps} configuration. The \gls{rtk-gps} solution will be calculated with the open source program rtklib, which will communicate with a task in DUNE. The solution from rtklib will be compared against the solution from Piksi, and a post processed solution from rtklib. The result from the experiment will be used in the discussion on how to perform a autonomous landing.

\section{System layout}
The current state of the system is that it consist of two parts that has not yet been integrated with each other. The part were which is the main focus of this thesis is the positioning part. The plan is to use \gls{rtk-gps} for positioning estimation. The second part is the path and control path. Currently there has only been done simulation of the guidance system. It shows promising results, but is yet to be integrated with \gls{rtk-gps}. For more information on the subject please read the MSc thesis by \citep{Froelich}. Figure \ref{figure:SystemOverview} gives a overview on how the system may look like. Ideas on how the intergration can be acheived will be given in the future work section. 

\begin{figure}[H]
	\centering
		\includegraphics[width=0.7\textwidth]{figs/SystemOverview.png}
		\caption{The structure of the autonomos landing system}
		\label{figure:SystemOverview}
\end{figure}
\section{Layout of thesis}
This section is currently not up to date

Chapter 1 Intro

Chapter 2 Theory about coordinate system

Chapter 3 Theory about GNSS system with focus on the GPS

Chapter 4 Hardware and software

Chapter 5 Test and result

Chapter 6 Conclusion, discussion and further work


%EVERYTHING BELLOW BEFORE CHAPTER 2 IS NOT PART OF THE THESIS.
%\begin{verbatim}
%\begin{eqnarray}\label{eq1}
%F = m \times a
%\end{eqnarray}
%\end{verbatim}
%
%\noindent This will produce
%
%\begin{eqnarray}\label{eq1}
%F = m \times a
%\end{eqnarray}
%
%\noindent To refer to the equation
%
%\begin{verbatim}
%\eqref{eq1}
%\end{verbatim}
%
%\noindent This will produce \eqref{eq1}.
%
%
%\section{Figures}
%To create a figure
%
%\begin{verbatim}
%\begin{figure}[h!]
%  \centering
%    \includegraphics[width=0.5\textwidth]{fig/pikachu}
%  \caption{Pikachu.}
%\label{fig1}
%\end{figure}
%\end{verbatim}
%
%
%
%\noindent To refer to the figure
%
%\begin{verbatim}
%\textbf{Fig. \ref{fig1}}
%\end{verbatim}
%
%\noindent This will produce \textbf{Fig. \ref{fig1}}
%
%\section{References}
%
%To cite references
%
%\begin{verbatim}
%\cite{1,2,3}
%\end{verbatim}
%or
%\begin{verbatim}
%\citep{1,2,3}
%\end{verbatim}
%
%\noindent This will produce: \cite{1,2,3} or \citep{1,2,3}, respectively.
%
%\section{Tables}
%
%To creat a table
%
%\begin{verbatim}
%\begin{table}[!h]
%\begin{center}
%    \begin{tabular}{ | l | l | l | l |}
%    \hline
%    \textbf{No.} & \textbf{Data 1} & \textbf{Data 2} \\ \hline
%     1 & a1 & b1 \\ \hline
%     2 & a2 & b2 \\ \hline
%    \end{tabular}
%\end{center}
%\caption{Table 1.}
%\label{Tab1}
%\end{table}
%\end{verbatim}
%
%\noindent This will produce
%
%\begin{table}[!h]
%\begin{center}
%    \begin{tabular}{ | l | l | l | l |}
%    \hline
%    \textbf{No.} & \textbf{Data 1} & \textbf{Data 2} \\ \hline
%     1 & a1 & b1 \\ \hline
%     2 & a2 & b2 \\ \hline
%    \end{tabular}
%\end{center}
%\caption{Table 1.}
%\label{Tab1}
%\end{table}
%
%\noindent To refer to the table
%
%\begin{verbatim}
%\textbf{Table. \ref{Tab1}}
%\end{verbatim}
%
%\noindent This will produce \textbf{Table. \ref{Tab1}}.

\cleardoublepage